\documentclass[11pt]{article}
\usepackage[utf8]{inputenc}
\usepackage[margin=1in]{geometry}
\usepackage{amsmath,amssymb,amsfonts}
\usepackage{graphicx}
\usepackage[dvipsnames]{xcolor}
\usepackage{authblk}
\usepackage[hidelinks]{hyperref}
\usepackage{natbib}
\usepackage{booktabs}
\usepackage{caption}
\usepackage{titlesec}
\usepackage{microtype}
\captionsetup{font=small,labelfont=bf}
\titleformat{\section}{\large\bfseries}{\thesection}{0.6em}{}
\titleformat{\subsection}{\normalsize\bfseries}{\thesubsection}{0.6em}{}
\renewcommand{\arraystretch}{1.15}

% ---- result macros (from results/summary.json) ----
\newcommand{\panrlat}{0.48}     % NB-VAE latent rare-type F1 (pancreas)
\newcommand{\panrpca}{0.00}     % PCA+Leiden rare-type F1 (pancreas, 0.002)
\newcommand{\markovari}{0.91}   % IGMC-Markov ARI (pancreas, no knob)
\newcommand{\pcaari}{0.47}      % PCA+Leiden ARI (pancreas)
\newcommand{\euclari}{0.76}     % NB-VAE+Euclidean+Leiden ARI (pancreas)
\newcommand{\covnominal}{90}
\newcommand{\covachieved}{90.0} % pancreas
\newcommand{\kappamed}{104}
\newcommand{\discpbmc}{0.84}    % discreteness index, PBMC
\newcommand{\discpanc}{0.95}    % discreteness index, pancreas
\newcommand{\disccont}{0.63}    % discreteness index, paul15 (continuous)
\newcommand{\ambcont}{27}       % conformal ambiguous fraction, paul15 (%)
% ---- WS1/WS2 macros (multi-seed CIs, CITE-seq, benchmark, scale-nested conformal) ----
\newcommand{\snsimcorr}{0.98}   % scale-nested simultaneous coverage, union-bound allocation
\newcommand{\snsimunc}{0.90}    % scale-nested simultaneous coverage, uncorrected
\newcommand{\snrarefine}{0.39}  % pancreas worst-class coverage at finest scale (Bonferroni)
\newcommand{\snrarecoarse}{0.99}% ... rescued at coarse scales
\newcommand{\csari}{0.75}       % CITE-seq PCA ARI vs protein labels
\newcommand{\csmarkovari}{0.72} % CITE-seq IGMC-Markov ARI vs protein labels
\newcommand{\csdcfr}{0.88}      % CITE-seq DC F1, Fisher-Rao
\newcommand{\csdcpca}{0.78}     % CITE-seq DC F1, PCA
\newcommand{\csdceuc}{0.27}     % CITE-seq DC F1, NB-VAE Euclidean ruler (same latent)
\newcommand{\msnbari}{0.82}     % pancreas NB-VAE+Euclid ARI (multi-seed mean)
\newcommand{\mspcaarims}{0.50}  % pancreas PCA ARI (multi-seed mean)
\newcommand{\msnbrare}{0.41}    % pancreas NB-VAE+Euclid rare-F1 (multi-seed mean)
\newcommand{\mspcararems}{0.12} % pancreas PCA rare-F1 (multi-seed mean)
\newcommand{\mssegnb}{0.10}     % segerstolpe NB-VAE+Euclid rare-F1
\newcommand{\bnigmc}{0.716}     % benchmark overall scIB score, IGMC latent
\newcommand{\bnscanvi}{0.733}   % scANVI (label-supervised)
\newcommand{\bnharmony}{0.723}  % Harmony
\newcommand{\bnscvi}{0.707}     % scVI
\newcommand{\bnpca}{0.564}      % PCA
\newcommand{\bnigmcbatch}{0.756}% IGMC batch-correction subscore

\title{\vspace{-1.2cm}\bfseries The right ruler: information-geometric, multiscale and
topologically-aware clustering of single-cell transcriptomes with conformal confidence}

\author[1]{Gabriella \and Megan \and Emma \and Rajarshi}
\affil[1]{Carnegie Mellon University Pre-College Computational Biology Program}
\date{}

\begin{document}
\maketitle
\vspace{-0.8cm}

\begin{abstract}
\noindent Single-cell RNA sequencing measures the transcriptome of individual cells, and
clustering those measurements into cell types is the field's central inference. Yet the
dominant pipeline is a chain of mutually inconsistent heuristics: it embeds counts with a
generative model, then discards that model's geometry to measure cell--cell distance with a
plain Euclidean ruler; it detects communities at a single, hand-tuned resolution; it cannot
say whether a population is a discrete type or a point on a continuum; and it reports every
assignment without a shred of calibrated confidence. We show that a negative-binomial
variational autoencoder already carries the correct geometry---the \emph{Fisher--Rao pullback
metric} that its own decoder Jacobian and count likelihood induce on the latent space---and
that this single object can do all four inferential jobs at once. It defines a
statistically honest inter-cell distance; it is scanned across scales by Markov-stability
community detection, removing the resolution knob; it is probed by persistent homology to
separate discrete types from continua; and it supplies a geometry-native nonconformity score
whose conformal prediction sets give every cell a calibrated confidence set. On a nine-technology
human pancreatic-islet atlas the negative-binomial geometry recovers rare cell types that standard
pipelines miss almost entirely (rare-type F1 up to \panrlat{} vs.\ \panrpca{} for PCA+Leiden), and
Markov stability on the metric graph recovers the cell-type hierarchy with no resolution tuning at
all (adjusted Rand index \markovari{}, versus \pcaari{} for PCA+Leiden). On a
type-2-diabetes cohort, calibrating the confidence layer on healthy islets and testing it on
diabetic ones, coverage is maintained overall but degrades most for $\beta$- and $\gamma$-cells:
the conformal layer localizes the disease's transcriptional footprint to the islet's endocrine
compartment without ever being told which cells are diseased. We stress-test these claims: results
carry multi-seed confidence intervals; an orthogonal CITE-seq protein ground truth shows the
count-aware Fisher--Rao geometry best recovers a rare dendritic-cell population even where standard PCA
remains competitive on common types; and we prove a \emph{scale-nested} conformal predictor that is
simultaneously valid and provably coherent across the Markov-stability hierarchy of scales. The result
converts scRNA-seq clustering from a stack of disconnected heuristics into one coherent,
uncertainty-calibrated inference on the manifold of the count model.
\end{abstract}

\section{Introduction}
Every cell in the body carries the same genome but expresses a distinct set of RNA molecules,
and single-cell RNA sequencing (scRNA-seq) reads out that expression one cell at a time
\citep{macosko2015highly}. The measurement is a sparse matrix of integer molecule counts---tens
of thousands of genes by up to millions of cells---and the first question asked of it is almost
always the same: which cells are of the same type? Answering it computationally has become the
workhorse of modern cell biology, underpinning tissue atlases, developmental maps and disease
studies alike.

The community has converged on a remarkably uniform recipe: normalize and log-transform the
counts, reduce dimension (by PCA or a deep generative model), build a $k$-nearest-neighbour
graph under Euclidean distance, and partition it with the Leiden algorithm at a user-chosen
resolution \citep{traag2019leiden,wolf2018scanpy}. Each step is individually reasonable and
collectively load-bearing, and each hides a flaw that the next step inherits.

First, \textbf{the ruler is wrong}. Count noise in scRNA-seq is not homoscedastic: the variance
of a gene grows with its mean, following a negative binomial (NB) law $\mathrm{Var}=\mu+\mu^2/\theta$
\citep{svensson2020droplet,townes2019glmpca}. Euclidean distance on log-normalized counts treats a
unit difference as equally meaningful everywhere, which it is not; the transformation injects
spurious variance that survives into PCA and clustering and disproportionately corrupts the rare,
lowly-sampled populations that are often the biological point \citep{townes2019glmpca,ahlmann2023comparison}.
Strikingly, even deep generative models such as scVI \citep{lopez2018scvi}, which fit an explicit NB
likelihood, throw the induced geometry away and cluster their latent space with plain $L_2$.

Second, \textbf{the resolution is arbitrary}. Cell identity is genuinely hierarchical---type,
subtype, state---but Leiden returns a single partition whose granularity is set by a resolution
parameter with no principled value. Slightly different choices yield materially different numbers
of clusters, and there is no accepted way to report the hierarchy itself.

Third, \textbf{discrete and continuous are conflated}. Some populations are sharp, well-separated
types; others are stops along a continuous trajectory (differentiation, activation, disease
progression). Standard clustering forces a partition on both, and cannot say which is which.

Fourth, \textbf{there is no calibrated confidence}. A cluster label is reported as a hard fact even
for a cell that sits ambiguously between two types or belongs to a state absent from the reference.

We argue these are not four separate problems but one, and that their common solution is already
latent in the generative model. An NB-VAE learns a decoder $f:\mathbb{R}^d\!\to\!\Delta^{G-1}$ from a
low-dimensional latent space to gene-expression proportions. The family of distributions it
parameterizes is a statistical manifold, and Rao's theorem endows it with a canonical Riemannian
metric---the Fisher information \citep{rao1945information,amari2000methods}. Pulling that metric back
through the decoder \citep{arvanitidis2022pulling} yields a metric tensor on the latent space,
$M(z)=J_f(z)^{\!\top} I_{\mathrm{NB}}(f(z))\,J_f(z)$, that measures the \emph{statistical
distinguishability} of cells rather than their raw coordinate distance. Our contribution is to make
this single geometric object do all four jobs at once: it is the ruler, it is the operator scanned by
Markov stability across scales, it is the filtration probed by persistent homology, and it supplies
the nonconformity score for conformal prediction. We do not claim to be first to bring information
geometry to single cells---closed-form multinomial Fisher--Rao has been used directly in gene space
\citep{cai2025gaia}, and the pullback framework is general \citep{arvanitidis2022pulling}. We claim to
be the first to pull the NB Fisher--Rao metric through a generative decoder at scale for single cells,
and to unify it with multiscale community detection, a topological discrete/continuous test, and
calibrated coverage into one coherent inference. We back these claims with multi-seed confidence
intervals, an orthogonal CITE-seq protein ground truth, a multi-method integration benchmark, and a
new \emph{scale-nested} conformal predictor that is provably valid and coherent across the whole
hierarchy of scales.

\section{Results}

\subsection{One geometry from the count model}
We train a batch-conditional NB-VAE (Methods) and, for each cell, compute the exact Fisher--Rao
pullback metric $M(z)$. The Fisher information of the negative binomial with respect to its mean is,
by a short derivation (Methods), $I(\mu)=\theta/(\mu(\mu+\theta))=1/\mathrm{Var}(x)$: the correct ruler
downweights each gene by its count variance, precisely the heteroscedastic correction the data
demand (Fig.~1a). Because the latent dimension $d$ is small ($\sim\!10$) while the gene dimension $G$
is large, the decoder Jacobian is computed exactly and cheaply by forward-mode automatic
differentiation, and the per-cell metric is only a $d\times d$ matrix; no full Jacobian is ever
materialized, which resolves the scalability concern that has made pullback geometry impractical at
single-cell scale. The resulting metric is strongly non-Euclidean: its condition number has median
$\kappa\!\approx\!\kappamed$ (Fig.~1c), the Riemannian volume element varies by orders of magnitude
across the manifold (Fig.~1b), and, decomposed against ground-truth labels, the metric systematically
stretches between-type distances while contracting within-type distances (Fig.~1d)---the mechanism by
which it sharpens cell-type boundaries.

\subsection{The Fisher--Rao ruler rescues rare cell types}
We benchmark clustering on a human pancreatic-islet atlas of 16{,}382 cells spanning fourteen cell
types across nine sequencing technologies, whose labels derive from marker genes rather than from any
clustering pipeline and are therefore an unbiased ground truth. The headline is stark on rare
populations (Fig.~6): the count-aware NB-VAE geometry recovers the rarest third of cell
types---epsilon (32 cells), mast (42), macrophage (79)---with rare-type F1 up to \panrlat, whereas
PCA+Leiden and $k$-means on log-normalized data recover essentially none of them (F1 $\approx$
\panrpca). Standard normalization does not merely under-resolve rare cells; it erases them.
Because real-atlas labels can be argued circular, we confirm the effect on planted-ground-truth
simulations with rare populations down to $0.4\%$ abundance (Fig.~8): on labels we control exactly,
PCA+Leiden recovers essentially no rare cells (rare-type F1 $0.04$) while the full information-geometric
method recovers them at F1 $0.45$, and recall rises smoothly with abundance for the count-aware methods
but collapses for PCA below $\sim\!2\%$. Here the three-ruler ablation on a shared latent is
unambiguous: the negative-binomial Fisher--Rao ruler is best on both cluster agreement (ARI $0.95$)
and rare-type recovery, ahead of the plain decoder pullback $J^{\!\top}J$ and Euclidean $L_2$
(Fig.~8c). On real data the metric's advantage over Euclidean is smaller and setting-dependent, and,
revealingly, the raw decoder pullback without the Fisher weighting can be numerically ill-behaved,
fragmenting a nine-technology atlas into hundreds of clusters (Fig.~2c); the negative-binomial Fisher
information both sharpens rare structure and \emph{regularizes} the metric.

\subsection{Markov stability removes the resolution knob}
Rather than fix a resolution, we let a continuous-time random walk explore the Fisher--Rao graph and
keep the partitions that stay stable across Markov time $t$ (Fig.~3). As the walk lengthens, the
optimal number of communities descends through a staircase of plateaus, and the partitions on those
plateaus are exactly the ones on which independent optimizer seeds agree---the variation of
information across seeds collapses to zero (Fig.~3a,b). The block structure of the
partition-vs-partition VI matrix reveals the cell-type hierarchy directly (Fig.~3c), and an alluvial
diagram traces how fine clusters merge into coarse ones across scales (Fig.~3d). The result is not
merely tidy but accurate: on the islet atlas the most robust plateau reaches an adjusted Rand index
of \markovari{} against marker-gene ground truth---far above PCA+Leiden (\pcaari) and above even the
batch-corrected NB-VAE latent clustered with Euclidean Leiden (\euclari)---with no resolution
parameter set anywhere in the pipeline.

\subsection{Persistent homology tells types from continua}
Whether a population is a discrete type or a stop on a continuum is a topological question, and we
answer it with persistent homology computed under the Fisher--Rao distance (Fig.~4). On data of known
discrete structure (peripheral-blood immune types) the persistence barcode shows components that
merge only across a large density gap, and the fraction of ``transitional'' cells---those whose
Fisher--Rao neighbourhood spans more than one population---is low. On data of known continuous
structure (a myeloid differentiation trajectory) the same statistics invert: cells bridge smoothly
between populations and the transitional fraction more than doubles. The resulting discreteness index
cleanly orders the datasets---\disccont{} for the continuous trajectory versus \discpbmc{} for immune
blood and \discpanc{} for the islet atlas---and the per-population transitionality score localizes
\emph{which} cells sit on a continuum, turning a global dichotomy into a per-cell readout.

\subsection{Conformal prediction calibrates every call}
Finally, we attach a calibrated confidence set to every cell by split Mondrian conformal prediction
with a class-conditional guarantee (Methods). At target coverage $1-\alpha=\covnominal\%$ the achieved
coverage is \covachieved\% on the islet atlas (Fig.~5a), and the guarantee holds per class, not just on
average. The size of a cell's prediction set is itself informative: a singleton is a confident call, a
set of size $>1$ marks an ambiguous cell between types, and an empty set flags a cell that conforms to
no known type. The set sizes track biology: on discrete blood and islet data almost every cell is a
confident singleton ($<2\%$ ambiguous), whereas on the continuous differentiation trajectory the
ambiguous fraction rises to \ambcont\%, precisely because cells there genuinely lie between types. And
within a dataset, ambiguity concentrates where topology predicts---prediction-set size rises with a
population's transitional fraction (Fig.~5d)---so the four modules agree with one another because they
are computed from the same geometry.

\subsection{Application: rare islet types and localizing type-2 diabetes}
We apply the full pipeline to human islets in health and type-2 diabetes (T2D), the setting where
resolving rare populations and honest confidence matter clinically. The method resolves every islet
type across both conditions (Fig.~7a,b). We then use the confidence layer as a probe of disease: we
calibrate the conformal classifier on healthy islets alone and measure its coverage on the diabetic
cohort, a covariate shift from health to disease. Overall coverage is maintained near the $90\%$
target ($90.7\%$), showing the guarantee is robust to the shift---but the coverage is not uniform.
It degrades most for $\beta$-cells ($87.8\%$) and $\gamma$-cells ($81.0\%$), the two endocrine
populations whose transcriptomes are most perturbed in T2D, while exocrine ductal and acinar cells
stay at or above target (Fig.~7c,d). In other words, the calibrated confidence layer localizes the
transcriptional footprint of diabetes to the islet's endocrine compartment---the $\beta$-cell above
all---without ever being shown a disease label. We emphasize what we do \emph{not} find: at the
resolution of this cohort the diabetic $\beta$-cells largely retain their identity (the classifier's
$\beta$-probability is essentially unchanged), so we make no claim of wholesale dedifferentiation
\citep{talchai2012dedifferentiation,weng2023hnf1a}; the honest, reproducible signal is a
compartment-specific loss of coverage that our uncertainty quantification surfaces automatically.

\subsection{Interdisciplinary geometry: optimal transport and the diffusion spectrum}
Because the metric is a genuine geometric object, two further graduate-level lenses apply directly.
First, \emph{optimal transport} (the Monge--Kantorovich / Wasserstein theory that underlies
Waddington-OT in single cell \citep{schiebinger2019}) lets us measure disease as displacement: for
each islet cell type we compute the squared Wasserstein-2 distance between healthy and diabetic cells
in latent space against a within-healthy null (Fig.~10). $\beta$-cells (insulin) and $\alpha$-cells
(glucagon) are displaced most, and a $250$-fold bootstrap shows the $\beta$-cell displacement is
$2.7\times$ the null ($p<0.004$) with cleanly separated distributions---an independent,
transport-theoretic confirmation of the conformal coverage result, and a rare instance of two
different mathematical frameworks agreeing that diabetes concentrates in the insulin cell. Second,
\emph{spectral graph theory and random-matrix theory}: the eigenspectrum of the diffusion operator
exhibits a spectral gap whose location counts the metastable states of the manifold and coincides with
the Markov-stability plateaus (Fig.~11), while a Marchenko--Pastur bound separates the signal principal
components from the noise bulk. Distance, scale, transport and spectrum are all facets of the same
geometry.

\subsection{Cross-dataset synthesis}
Read together (Fig.~9), the four datasets tell one story. Markov-stability clustering is the most
accurate method on every fair benchmark; the count-aware latent recovers rare populations on all of
them; conformal coverage sits at its target throughout; and the topological discreteness index and
the conformal ambiguous fraction move in lockstep---both flag the myeloid continuum and both call
the immune and islet data discrete. An ablation confirms that the no-knob Markov result is helped by,
and never hurt by, building the graph on the Fisher--Rao metric rather than the Euclidean one
(Fig.~9d). Four readouts, computed from one geometry, agree.

\subsection{The wins survive multi-seed retraining}
Because the metric is read off a trained VAE, a fair critic asks whether the results are an artifact of
one lucky initialization. We retrain the NB-VAE from scratch across seeds and recompute every metric,
reporting mean and $95\%$ bootstrap confidence intervals (Fig.~12). The picture is honestly mixed and,
where it matters, decisive. On the fair pancreas benchmark the count-aware latent's advantage is
robust: NB-VAE+Euclidean reaches ARI $\msnbari$ (95\% CI $[0.78,0.86]$) versus $\mspcaarims$
$[0.48,0.52]$ for PCA+Leiden---intervals that do not overlap---and rare-type F1 $\msnbrare$ versus
$\mspcararems$. On the Segerst{\o}lpe islets PCA+Leiden recovers \emph{zero} rare-type F1 in every
single seed ($0.00\,[0.00,0.00]$), while the count-aware latent recovers $\mssegnb$. Conversely, on the
circular-label PBMC3k data PCA is both higher and more stable (ARI $0.87\,[0.86,0.87]$) than the
count-aware rulers---consistent with those labels having been produced by a PCA pipeline. Conformal
coverage is stable at its $90\%$ target across all seeds and datasets (Fig.~12d). The rigor cuts both
ways, and we report it that way: the count-aware geometry's gains are real and seed-robust precisely on
the rare-population and batch-integration problems it was designed for.

\subsection{An orthogonal ground truth from CITE-seq protein}
A standing objection to any scRNA-seq clustering benchmark is circularity: the ``ground-truth''
labels are themselves the output of an RNA pipeline (often PCA + Louvain), so a method that clusters
like that pipeline is rewarded regardless of biological truth. We break the circle with a modality
that never touches the RNA counts. On a 10x CITE-seq PBMC dataset ($4{,}193$ cells) we define cell
types from $32$-antibody surface-protein measurements alone---centred-log-ratio normalisation,
Leiden clustering on protein, and canonical marker gating (CD3/CD4/CD8, CD19/CD20, CD14/CD16,
CD56, HLA-DR)---yielding six protein-defined populations including a rare ($4.5\%$) dendritic-cell
compartment (Fig.~14a,b). We then cluster the \emph{RNA} counts with every method and score against
these orthogonal labels. The honest result has two parts. First, against orthogonal truth PCA+Leiden
remains strong on the common types (ARI $\csari$, above IGMC's $\csmarkovari$)---so PCA's advantage
on discrete PBMC data is real, not merely a circularity artefact, and we say so. Second, and
consistent with our central thesis, the Fisher--Rao geometry best recovers the \emph{rare}
dendritic-cell type (DC F1 $\csdcfr$ vs.\ $\csdcpca$ for PCA), and---critically---the very same
count-aware latent scored with a plain Euclidean ruler collapses on this population (DC F1 $\csdceuc$;
Fig.~14c,d). The count-aware latent plus the Fisher--Rao ruler is what rescues the rare type, and it
does so measured against a ground truth no RNA method could have inflated.

\subsection{Competitive with established integration methods}
Reviewers rightly ask how the count-aware latent fares against the field's dedicated tools. We run the
scIB battery (bio-conservation and batch-removal, implemented transparently from primitives) on the
nine-technology pancreas atlas against PCA, ComBat, Harmony \citep{korsunsky2019harmony}, scVI and the
label-supervised scANVI \citep{lopez2018scvi,gayoso2022scvitools}, all at ten latent dimensions
(Fig.~15). The NB-VAE latent is competitive with the best: an overall score of $\bnigmc$ places it
third of six---behind only scANVI ($\bnscanvi$), which is \emph{semi-supervised} and sees the cell-type
labels, and Harmony ($\bnharmony$), a method built solely for batch removal---and ahead of scVI
($\bnscvi$), ComBat, and PCA ($\bnpca$, worst by a wide margin). It achieves the second-best batch
mixing of any method ($\bnigmcbatch$). We flag its one clear weakness honestly: the batch-isolated-label
score is low ($0.003$), because the batch-conditional decoder removes technical structure aggressively
and can over-mix a cell type confined to a single technology. The headline is not that the geometry
wins a leaderboard---it is that a latent built for a \emph{different} purpose (a statistically honest
metric) is already on par with purpose-built integrators, while additionally supplying the scale
selection, topology and calibrated confidence they do not.

\subsection{Scale-nested conformal prediction across the hierarchy}
A clustering is only meaningful \emph{at a scale}, so a confidence set should come with a scale too.
The Markov-stability scan already yields cell communities from fine to coarse; tracking which cell
types are walked together across Markov time induces a nested \emph{taxonomy} (Fig.~13a, biologically
sensible: endocrine, exocrine, immune and stromal super-types emerge). We attach to every cell a whole
family of conformal prediction sets, one per taxonomic level, and establish three properties (Methods;
full proofs in the supplement). (i)~Each level is a valid split-conformal predictor---standard
\citep{vovk2005algorithmic,sadinle2019least}. (ii)~Allocating the miscoverage budget across levels
($\sum_\ell\alpha_\ell=\alpha$) gives a \emph{simultaneous} family-wise guarantee by a union
bound---also standard \citep{baheri2025multiscale}; empirically the guarantee holds with margin
(simultaneous coverage $\approx\snsimcorr$ at target $0.90$), whereas running every level at the full
$\alpha$ sits at the target with none ($\approx\snsimunc$; Fig.~13c). (iii)~Our new contribution is a
construction that makes the family \emph{coherent}---a confident fine-scale call can never contradict
its coarse-scale super-type---by an isotonic up-the-tree projection of the per-scale thresholds that
we prove preserves the per-scale coverage (Theorem, Methods). The projection raises coherence to
$100\%$ from $90$--$97\%$ (Fig.~13d), eliminating the $3$--$10\%$ of cells that would otherwise carry a
fine call orphaned from its coarse super-type, and it sits between isotonic-calibration
\citep{vanderlaan2024selfcalibrating} and nested-set \citep{gupta2022nested} conformal lines, distinct
from fixed-taxonomy conformal cell typing \citep{corbetta2024conformal,lopezdecastro2025conformal,ding2023classcond}.
The hierarchy is not decoration: a rare pancreatic type covered at only $\snrarefine$ at the finest
scale (too few calibration cells to certify class-conditionally) is rescued to $\snrarecoarse$ once it
is pooled into a well-populated super-type (Fig.~13b,e), and the residual ambiguity that survives
coarsening concentrates on the myeloid continuum (Fig.~13f)---the same population the topology flags.

\section{Discussion}
We have shown that the geometry a negative-binomial generative model induces on its latent space is
not a curiosity to be discarded but the organizing object of single-cell clustering. Reading distance,
scale, shape and confidence off the same Fisher--Rao pullback metric replaces four inconsistent
heuristics with one coherent inference, and does so with concrete gains: rare cell types that standard
pipelines erase are recovered, the resolution knob is eliminated, discrete and continuous structure are
distinguished, and every assignment carries a finite-sample coverage guarantee.

Our study has clear limits, which also chart the path forward. We are deliberately careful about what
the orthogonal benchmark shows: against protein-defined labels PCA remains competitive on common PBMC
types, so the count-aware geometry's advantage is specifically a rare-population and continuum effect,
not a uniform improvement. The scale-nested predictor likewise trades informativeness for guarantee---the
provable family-wise allocation is conservative, and its union bound is loose under the strong positive
dependence that nesting induces; a tighter valid joint calibration is open. The induced metric depends
on the fitted decoder; marginalizing over a decoder ensemble \citep{syrota2024decoder} would render the
geometry identifiable and is a natural next step. Our evaluation spans four datasets up to $\sim\!16$k cells;
the graph-geodesic and low-rank spectral machinery we introduce is designed to scale, but a
demonstration on million-cell atlases remains. The benefit of the Fisher--Rao ruler over a plain
decoder-pullback is clearest where batch structure is controlled, underscoring that geometry and batch
correction must be co-designed. Finally, the diabetic-$\beta$-cell readout is a compartment-specific
coverage drop, not a dedifferentiation claim; connecting that statistical signal to specific
transcriptional programs, and validating it in an independent cohort, is future work. None of these
caveats undercuts the central point: single-cell clustering has been done with the wrong ruler, and a
principled one---together with the topology and calibrated confidence it unlocks---was inside the
generative model all along.

\section{Methods}

\paragraph{Negative-binomial VAE.} We model UMI counts $x_{ig}$ for cell $i$, gene $g$ as
$x_{ig}\sim\mathrm{NB}(\mu_{ig},\theta_g)$ with mean $\mu_{ig}=\ell_i\,\rho_{ig}$, library size
$\ell_i=\sum_g x_{ig}$, gene proportions $\rho_i=\mathrm{softmax}(f_{\text{dec}}([z_i,b_i]))$ on the
simplex, and gene-specific inverse dispersion $\theta_g=\mathrm{softplus}(\cdot)$, so that
$\mathrm{Var}(x_{ig})=\mu_{ig}+\mu_{ig}^2/\theta_g$. An encoder $q(z_i\mid x_i,b_i)=\mathcal{N}(\mu_z,\mathrm{diag}\,\sigma_z^2)$
and the decoder are two-layer perceptrons with LayerNorm (no BatchNorm, so the decoder is compatible
with functional automatic differentiation); $b_i$ is a one-hot batch covariate supplied to both
encoder and decoder, which pushes technical variation out of $z$. We optimize the $\beta$-VAE ELBO
$\sum_i \mathbb{E}_q[-\log p(x_i\mid z_i,b_i)]+\beta\,\mathrm{KL}(q\|\,\mathcal{N}(0,I))$ with
$\beta$ warmed from 0 to 1, Adam, gradient clipping, $d=10$.

\paragraph{Fisher--Rao pullback metric.} The Fisher information of $\mathrm{NB}(\mu,\theta)$ with
respect to $\mu$ is $I(\mu)=-\mathbb{E}\,\partial_\mu^2\log p = 1/\mu - 1/(\mu+\theta) =
\theta/(\mu(\mu+\theta)) = 1/\mathrm{Var}(x)$. Treating $\{\mathrm{NB}(\mu(z),\theta)\}$ as a
statistical manifold and pulling the Fisher metric back through the decoder gives, on the latent space,
\[
  M(z)=\sum_g I(\mu_g)\,(\partial_z \mu_g)(\partial_z \mu_g)^{\!\top}
      =\ell^2\, J_\rho(z)^{\!\top}\,\mathrm{diag}\!\big(\theta_g/(\mu_g(\mu_g+\theta_g))\big)\,J_\rho(z),
\]
where $J_\rho=\partial\rho/\partial z\in\mathbb{R}^{G\times d}$. Because $d\ll G$ we compute $J_\rho$ by
forward-mode AD (\texttt{jacfwd}) vectorized over cells (\texttt{vmap}), forming $M(z)$ as a $d\times d$
matrix per cell without materializing any $G\times G$ object; the whole atlas takes seconds. The local
squared Fisher--Rao distance between neighbours uses the symmetrized midpoint metric
$\delta^2(i,j)=(z_i-z_j)^{\!\top}\tfrac12(M_i+M_j)(z_i-z_j)$, a second-order-accurate geodesic
surrogate. The ablation metric replaces $I(\mu)$ by $1$, giving the plain decoder pullback
$M=\ell^2 J_\rho^{\!\top}J_\rho$.

\paragraph{Metric graph.} We build a $k$-nearest-neighbour graph by candidate-and-rerank: fetch
$k_{\text{cand}}\!=\!60$ Euclidean candidates per cell (cheap), rescore them under $\delta^2$, keep the
$k=15$ statistically nearest, weight edges with an adaptive Gaussian kernel and symmetrize. Cost is
$O(N k_{\text{cand}} d^2)$, avoiding any $N^2$ blow-up.

\paragraph{Markov-stability multiscale clustering.} On the metric graph with adjacency $A$, degrees
$D$, stationary $\pi_i=d_i/2m$, we consider the continuous-time random walk $P(t)=\exp(-t(I-D^{-1}A))$.
The clustered stability of a partition $H$ at Markov time $t$ is
$r(t,H)=\mathrm{tr}\!\big[H^{\!\top}(\Pi P(t)-\pi\pi^{\!\top})H\big]$. Because the walk is reversible,
$\Pi P(t)$ is symmetric and non-negative, and the graph with adjacency $2m\,\Pi P(t)$ has the same
degrees as $A$; hence $r(t,H)$ equals the RB-modularity of that flow graph at resolution 1, which
Leiden optimizes directly. We build $\Pi P(t)=\tfrac1{2m}W\mathrm{diag}(e^{-t\Lambda})W^{\!\top}$ from
the eigendecomposition of the symmetric normalized Laplacian ($W=D^{1/2}U$), sparsify to the top edges
per row, and optimize with several seeds over a log grid of $t$. Robust scales are Markov times where
the number of communities forms a plateau and the across-seed variation of information
\citep{meila2007comparing} dips; large atlases are stratified-subsampled for the scan to preserve rare
types.

\paragraph{Persistent homology.} For a population we compute Vietoris--Rips persistence (Ripser) under
the Fisher--Rao distance. The dataset ``continuity index'' is the fraction of cells whose $k$
Fisher--Rao neighbours span more than one population (transitional cells; the $H_0$ bridges of the
complex); its complement is the discreteness index. Global $H_0$/$H_1$ diagrams and the $H_0$ death-gap
statistic summarize manifold shape.

\paragraph{Conformal prediction.} We use split conformal with a Mondrian (class-conditional) calibration
and the least-ambiguous set-valued score $s(x,y)=1-\hat p(y\mid x)$ from a multinomial classifier on
$z$ \citep{sadinle2019least,angelopoulos2023gentle}. The per-class threshold is the finite-sample
$\lceil(n_y+1)(1-\alpha)\rceil/n_y$ quantile of calibration scores, giving prediction sets
$C(x)=\{y:\hat p(y\mid x)\ge 1-\hat q_y\}$ with $\Pr(y\in C(x)\mid Y=y)\ge 1-\alpha$ per class. Sets of
size $>1$ mark ambiguity; empty sets mark novelty. For the disease analysis we calibrate on healthy
cells and evaluate on the diabetic cohort.

\paragraph{Scale-nested conformal prediction.} A taxonomy has levels $\ell=1,\dots,L$ (finest to
coarsest) with coarsening maps $\pi_\ell:\mathcal Y\to\mathcal G_\ell$ and parent maps
$\rho_\ell:\mathcal G_\ell\to\mathcal G_{\ell+1}$ satisfying $\pi_{\ell+1}=\rho_\ell\circ\pi_\ell$;
we build it by average-linkage on the integrated Markov co-clustering distance
$D_{ab}=1-\frac1T\sum_t\sum_c P(a\in c\mid t)P(b\in c\mid t)$ between cell types $a,b$. Super-class
probabilities aggregate a base classifier, $p_\ell(g\mid x)=\sum_{y:\pi_\ell(y)=g}\hat p(y\mid x)$;
the level-$\ell$ set is $C_\ell(x)=\{g:p_\ell(g\mid x)\ge 1-q_\ell\}$ with $q_\ell$ the split-conformal
quantile at level $\alpha_\ell$. \emph{(1)} Each level is marginally valid,
$\Pr(\pi_\ell(Y)\in C_\ell(X))\ge1-\alpha_\ell$ \citep{vovk2005algorithmic}. \emph{(2)} With
$\sum_\ell\alpha_\ell=\alpha$, a union bound gives simultaneous validity
$\Pr(\exists\ell:\pi_\ell(Y)\notin C_\ell(X))\le\alpha$ \citep{baheri2025multiscale}. \emph{(3, new)}
Define projected thresholds by the fine$\to$coarse pass $\tilde q_1=q_1$,
$\tilde q_{\ell+1}=\max(q_{\ell+1},\tilde q_\ell)$. Because $\{s\le q\}\subseteq\{s\le\tilde q\}$ for
$\tilde q\ge q$, the projection only enlarges sets and preserves (1)--(2); and because
$p_{\ell+1}(\rho_\ell(g)\mid x)\ge p_\ell(g\mid x)$ (nesting) with $\tilde q_{\ell+1}\ge\tilde q_\ell$,
$g\in C_\ell(x)$ implies $\rho_\ell(g)\in C_{\ell+1}(x)$: the family is \emph{coherent}. We validate on
held-out cells over $25$ splits; the per-cell resolvable scale is the finest level with a singleton
set, cross-fitted out-of-fold. Full statements and proofs are in the supplement.

\paragraph{CITE-seq orthogonal ground truth.} From the 10x $5$k-PBMC TotalSeq-B dataset
\citep{stoeckius2017citeseq} we split gene-expression from the $32$ antibody-capture features. Protein
is CLR-normalised (per cell, isotype controls excluded), PCA/Leiden-clustered, and each cluster is
gated to a canonical PBMC lineage by fixed marker rules. RNA counts are then processed and clustered
identically to the other datasets and scored against these protein-only labels.

\paragraph{Multi-method benchmark and multi-seed CIs.} We compare embeddings from PCA, ComBat, Harmony
\citep{korsunsky2019harmony}, scVI and scANVI \citep{lopez2018scvi,gayoso2022scvitools}, and our
count-aware NB-VAE latent, all at $10$ latent dimensions, using a from-scratch scIB battery
\citep{luecken2022benchmarking} (NMI, ARI, cell-type ASW, isolated-label F1, cLISI; batch ASW, graph
connectivity, iLISI; overall $=0.6\,\mathrm{bio}+0.4\,\mathrm{batch}$). To quantify seed sensitivity we
retrain the NB-VAE from six random seeds (five on the larger pancreas atlas) and report mean and $95\%$
bootstrap confidence intervals for every headline metric.

\paragraph{Optimal transport and spectral geometry.} Disease displacement is the squared
Wasserstein-2 distance $W_2^2(\mu_{\text{healthy}},\mu_{\text{T2D}})$ between per-type empirical
distributions in latent space (exact OT via the Python Optimal Transport library), reported against a
within-healthy split null and a $250$-fold paired bootstrap. The diffusion analysis eigendecomposes the
symmetric normalised Laplacian of the latent kNN graph; the largest gap in the low end of the spectrum
estimates the number of metastable states, and relaxation times $1/\lambda_k$ are compared to the
Markov-stability plateaus. The Marchenko--Pastur edge $\lambda_+=(1+\sqrt{p/n})^2$ on the
standardised expression matrix separates signal principal components from noise.

\paragraph{Data and evaluation.} We use PBMC 3k (10x, 2{,}638 cells), the Paul et al.\ myeloid
differentiation set (2{,}730), a nine-technology human pancreatic islet atlas (16{,}382 cells, 14
types) and the Segerst{\o}lpe healthy-vs-T2D islet cohort (2{,}394 cells), all open access with integer
counts. Baselines are PCA+Leiden, $k$-means on the latent, NB-VAE latent with Euclidean Leiden, and the
decoder-pullback ablation. Metrics are ARI, NMI, silhouette and a rare-type F1 that averages
Hungarian-matched per-type F1 over the rarest tertile of cell types. Code, trained models and figures
are released with the paper.

\clearpage
\section*{Appendix: Scale-nested conformal prediction --- statements and proofs}
\setlength{\parskip}{2pt}
\textbf{Setup.} Taxonomy levels $\ell=1,\dots,L$ (finest to coarsest), coarsening maps
$\pi_\ell:\mathcal Y\to\mathcal G_\ell$ and parent maps $\rho_\ell:\mathcal G_\ell\to\mathcal G_{\ell+1}$
with $\pi_{\ell+1}=\rho_\ell\circ\pi_\ell$. A base classifier gives $\hat p(y\mid x)$; super-class
scores aggregate it, $p_\ell(g\mid x)=\sum_{y:\pi_\ell(y)=g}\hat p(y\mid x)$, with LAC nonconformity
$s_\ell(x,g)=1-p_\ell(g\mid x)$. Given exchangeable calibration $(X_i,Y_i)_{i=1}^n$, the marginal
threshold $q_\ell$ is the $\lceil(n+1)(1-\alpha_\ell)\rceil$-th smallest of $\{s_\ell(X_i,\pi_\ell(Y_i))\}$
(and $q_\ell=1$ if that rank exceeds $n$); the set is $C_\ell(x)=\{g:p_\ell(g\mid x)\ge 1-q_\ell\}$.

\textbf{Proposition 1 (per-level validity; standard).} For each $\ell$,
$\Pr(\pi_\ell(Y_{n+1})\in C_\ell(X_{n+1}))\ge 1-\alpha_\ell$.
\emph{Proof.} The event is $\{s_\ell(X_{n+1},\pi_\ell(Y_{n+1}))\le q_\ell\}$; by exchangeability the
rank of the test score among the $n+1$ scores is uniform on $\{1,\dots,n+1\}$, so the probability is at
least $\lceil(n+1)(1-\alpha_\ell)\rceil/(n+1)\ge 1-\alpha_\ell$ \citep{vovk2005algorithmic,sadinle2019least}.
$\square$

\textbf{Proposition 2 (simultaneous validity; standard union bound).} If $\sum_\ell\alpha_\ell=\alpha$
then $\Pr(\exists\,\ell:\pi_\ell(Y_{n+1})\notin C_\ell(X_{n+1}))\le\alpha$.
\emph{Proof.} Boole's inequality over the $L$ miscoverage events, each of probability $\le\alpha_\ell$ by
Prop.~1 \citep{baheri2025multiscale}. $\square$

\textbf{Lemma A (monotone aggregation).} For $g\in\mathcal G_\ell$ with parent $h=\rho_\ell(g)$,
$p_{\ell+1}(h\mid x)\ge p_\ell(g\mid x)$.
\emph{Proof.} Nesting gives $\{y:\pi_\ell(y)=g\}\subseteq\{y:\pi_{\ell+1}(y)=h\}$; probabilities are
nonnegative, so the larger index set has the larger sum. $\square$

\textbf{Lemma B (monotonization preserves coverage).} Replacing $q$ by any $\tilde q\ge q$ enlarges
$C_\ell$ ($\{s\le q\}\subseteq\{s\le\tilde q\}$) and cannot decrease the Prop.~1 coverage. $\square$

\textbf{Theorem 3 (coherent construction preserving validity; new).} Define
$\tilde q_1=q_1$ and $\tilde q_{\ell+1}=\max(q_{\ell+1},\tilde q_\ell)$ (fine$\to$coarse). Then
$C_\ell(x)=\{g:p_\ell(g\mid x)\ge 1-\tilde q_\ell\}$ satisfies Props.~1--2 \emph{and} is coherent:
$\rho_\ell(C_\ell(x))\subseteq C_{\ell+1}(x)$ for all $x,\ell$.
\emph{Proof.} Validity: $\tilde q_\ell\ge q_\ell$, so Lemma~B preserves each per-level bound and hence
the union bound. Coherence: if $g\in C_\ell(x)$ then, with $h=\rho_\ell(g)$,
$p_{\ell+1}(h\mid x)\ge p_\ell(g\mid x)\ge 1-\tilde q_\ell\ge 1-\tilde q_{\ell+1}$, using Lemma~A and
$\tilde q_{\ell+1}\ge\tilde q_\ell$; thus $h\in C_{\ell+1}(x)$. $\square$

These are finite-sample and distribution-free (exchangeability only). Props.~1--2 are standard; Thm.~3
--- the isotonic-threshold construction that makes a taxonomic family coherent while provably preserving
per-scale coverage --- is the contribution. It relates to isotonic-in-conformal calibration
\citep{vanderlaan2024selfcalibrating} and nested conformal sets \citep{gupta2022nested}, and differs
from fixed-taxonomy conformal cell typing \citep{corbetta2024conformal,ding2023classcond}.

\clearpage
\section*{Figures}
\newcommand{\figblock}[3]{\begin{figure}[p]\centering
\includegraphics[width=\linewidth]{../figures/#1}\caption{#2}\label{#3}\end{figure}}

\figblock{fig1_geometry_pbmc3k.pdf}{\textbf{Information geometry of the negative-binomial count
manifold.} (a) Gene variance grows with the mean, following the NB law; the Fisher information
$I(\mu)=1/\mathrm{Var}$ (inset) is the correct per-gene weight. (b) The Fisher--Rao metric field in
a latent projection, with distance ellipses and the log-volume (magnification) as background.
(c) The metric is strongly anisotropic (median condition number $\kappa\!\approx\!100$). (d) Under
the metric, between-type distances are stretched relative to within-type distances.}{fig:geom}

\figblock{fig2_clustering_pbmc3k.pdf}{\textbf{The right ruler resolves cell types.} Euclidean (a)
vs.\ Fisher--Rao (b) embeddings. (c) Ablation on a shared latent across three rulers---plain $L_2$,
decoder pullback $J^{\!\top}J$, and NB Fisher--Rao---isolating the count geometry's contribution.
(d) Per-type F1 improves for most populations under the Fisher--Rao ruler.}{fig:clust}

\figblock{fig3_markov_pbmc3k.pdf}{\textbf{Markov-stability multiscale clustering.} (a) Number of
communities and stability across Markov time, with robust plateaus shaded. (b) Across-seed variation
of information dips at robust scales. (c) The partition-vs-partition VI matrix reveals multiscale
block structure. (d) An alluvial diagram of the cluster hierarchy across scales.}{fig:markov}

\figblock{fig4_topology.pdf}{\textbf{Persistent homology separates discrete types from continua.}
Persistence diagrams for discrete (a) and continuous (b) data under the Fisher--Rao distance; the
transitional-cell map (c); and per-population transitionality across regimes (d).}{fig:topo}

\figblock{fig5_conformal_pancreas.pdf}{\textbf{Conformal prediction.} (a) Achieved coverage matches
the target, per class. (b) Prediction-set composition (confident / ambiguous / novel) across
$\alpha$. (c) The ambiguity map. (d) Prediction-set size rises with a population's topological
transitionality.}{fig:conf}

\figblock{fig6_pancreas_pancreas.pdf}{\textbf{Human pancreatic islet atlas.} (a) The Fisher--Rao
embedding with rare types labeled. (b) Batch mixing across nine technologies. (c) The fair-benchmark
method comparison. (d) Per-type recovery of rare islet populations.}{fig:panc}

\figblock{fig7_disease_segerstolpe.pdf}{\textbf{Type-2-diabetes stress test.} Islets by cell type (a)
and disease (b). (c) Diabetic $\beta$-cells are significantly more topologically transitional. (d)
Calibrated on healthy cells, the conformal layer assigns diabetic $\beta$-cells more ambiguous and
novel prediction sets.}{fig:dis}

\figblock{fig8_simulation.pdf}{\textbf{Planted-ground-truth simulation.} (a) Simulated data with
rare populations. (b) Per-type recall as a function of planted abundance. (c) The ruler ablation on
labels we control exactly. (d) Rare-type recovery across all methods.}{fig:sim}

\figblock{fig9_synthesis.pdf}{\textbf{Cross-dataset synthesis.} (a) Clustering accuracy (ARI) for
every method on every dataset. (b) Rare-population recovery, standard vs.\ count-aware. (c) Conformal
coverage, discreteness index and ambiguous fraction across datasets---confidence and shape agree.
(d) The Fisher--Rao graph helps, and never hurts, the no-resolution-knob Markov result.}{fig:synth}

\figblock{fig10_optimaltransport.pdf}{\textbf{Optimal transport of disease.} (a) Squared-$W_2$
displacement between healthy and diabetic cells per type, against a within-healthy null (ticks).
(b) Displacement above null: insulin ($\beta$) and glucagon ($\alpha$) cells move most. (c) The
$\beta$-cell transport map, healthy $\to$ T2D. (d) A $250$-fold bootstrap: the $\beta$-cell shift is
$2.7\times$ the null, $p<0.004$.}{fig:ot}

\figblock{fig11_spectral.pdf}{\textbf{Spectral \& diffusion geometry with a random-matrix null.}
(a) The diffusion eigenspectrum; the largest spectral gap counts metastable states. (b) A
Marchenko--Pastur bound separates signal principal components from the noise bulk. (c) Diffusion
relaxation times coincide with the Markov-stability plateaus. (d) The spectral signature differs by
tissue.}{fig:spec}

\figblock{fig14_citeseq_orthogonal.pdf}{\textbf{CITE-seq protein as an orthogonal ground truth.}
(a) The RNA Fisher--Rao embedding, coloured by cell types defined from surface protein alone
(orthogonal to any RNA pipeline). (b) Mean CLR protein markers per type---the labels are protein-driven.
(c) RNA clustering scored against the protein labels: PCA remains strong on common types, but the
count-aware methods lead on rare-type F1. (d) Recovery of the rare dendritic-cell type by method: the
Fisher--Rao ruler wins (F1 $\csdcfr$), and the same latent read with a Euclidean ruler fails
($\csdceuc$)---isolating the geometry's contribution on orthogonal truth.}{fig:cite}

\figblock{fig13_scale_nested_conformal.pdf}{\textbf{Scale-nested conformal prediction on the
Markov-stability hierarchy.} (a) The nested cell-type taxonomy from integrated Markov co-clustering
(pancreas). (b) Per-scale coverage: marginal coverage holds at every level, while a rare type
under-served at the finest scale ($\snrarefine$) is rescued as it is pooled into coarse super-types.
(c) Simultaneous (family-wise) coverage: the union-bound allocation buys a margin over the $0.90$
target; an uncorrected family sits at the target with none. (d) Coherence: the isotonic-projection
construction makes every cell coherent (fine calls never contradict their coarse super-type), versus
$90$--$97\%$ without it. (e) Rare-type rescue generalises across datasets. (f) Prediction sets
stay compact; the continuum (Paul15) carries the most residual ambiguity. Held-out, 25 splits.}{fig:snc}

\figblock{fig12_multiseed.pdf}{\textbf{Multi-seed confidence intervals.} Retraining the NB-VAE from
$10$ seeds ($6$ for the pancreas atlas), the headline results are not a lucky seed. (a)~Clustering
accuracy (ARI) and (b)~rare-type F1 with $95\%$ bootstrap CIs across methods and datasets; the
count-aware advantage separates from PCA. (c)~Per-seed spread on the pancreas atlas. (d)~Conformal
coverage sits at its $90\%$ target across seeds.}{fig:ms}

\figblock{fig15_benchmark.pdf}{\textbf{Benchmarking against established integration methods.} The
transparent scIB battery (bio-conservation and batch-correction) on the nine-technology pancreas atlas,
comparing PCA, ComBat, Harmony, scVI, scANVI and the count-aware NB-VAE latent. (a)~The
bio-conservation vs.\ batch-correction trade-off. (b)~Overall score. (c,d)~Metric breakdowns. Scanorama
is omitted (its dependency requires a compiler unavailable on our hardware).}{fig:bench}

\clearpage
\bibliographystyle{unsrtnat}
\bibliography{refs}
\end{document}
